%%%%%%%%%%%%%%%%%%%%%%%%%%%%%%%%%%%%%%%%%%%%%%%%%%%%%%%%%%%%%%%%%%%%%%%%%%
% JetBrains -- AI Lead, Python Tools (PyCharm). Tailored resume.
% Framing: AI-native developer tooling, coding agents, MCP, LLM evaluation/benchmarking.
% Fonts: Source Sans Pro. Accent: teal + purple. Compile: pdflatex (2 passes).
%%%%%%%%%%%%%%%%%%%%%%%%%%%%%%%%%%%%%%%%%%%%%%%%%%%%%%%%%%%%%%%%%%%%%%%%%%

\documentclass[10pt,a4paper]{article}

\usepackage[T1]{fontenc}
\usepackage[utf8]{inputenc}
\usepackage[default]{sourcesanspro}
\usepackage[a4paper,top=1.1cm,bottom=1.0cm,left=1.5cm,right=1.5cm]{geometry}
\usepackage{titlesec}
\usepackage{enumitem}
\usepackage{xcolor}
\usepackage{hyperref}
\usepackage{microtype}

% ---- Colours (teal primary, purple accent) ----
\definecolor{accent}{HTML}{117A8B}
\definecolor{company}{HTML}{6C2BB3}
\definecolor{muted}{HTML}{555555}

\hypersetup{colorlinks=true, urlcolor=accent, linkcolor=accent}

\setlength{\parindent}{0pt}
\linespread{1.0}
\pagestyle{empty}

\titleformat{\section}
  {\normalfont\large\bfseries\color{accent}\scshape}
  {}{0em}{}[{\color{accent!35}\titlerule[1pt]}]
\titlespacing*{\section}{0pt}{7pt}{4pt}

\newlist{tight}{itemize}{1}
\setlist[tight]{leftmargin=1.3em,itemsep=1pt,topsep=1pt,parsep=0pt,label=\textcolor{accent}{\small$\bullet$}}

\newcommand{\entry}[4]{%
  {\color{company}\bfseries #1}\hfill {\color{muted}#2}\par
  \vspace{1pt}
  {\itshape #3}\hfill {\color{muted}\itshape #4}\par
  \vspace{3pt}%
}

\newcommand{\sep}{\,\textcolor{accent!60}{\textbar}\,}

\begin{document}

%======================= HEADER =======================
\begin{center}
  {\fontsize{24}{26}\selectfont\bfseries Abhijeet Lokhande}\\[3pt]
  {\color{accent}\large AI Engineer: AI-Native Developer Tooling, Agents \& MCP}\\[5pt]
  {\small
    London, United Kingdom \sep
    \href{mailto:abhijeet_lokhande@proton.me}{abhijeet\_lokhande@proton.me} \sep
    07480801382 \\[2pt]
    \href{https://www.linkedin.com/in/ablds/}{LinkedIn} \sep
    \href{https://github.com/abhijeetscode}{GitHub} \sep
    \href{https://huggingface.co/ab-ai}{Hugging Face}
  }
\end{center}
\vspace{2pt}

%======================= SUMMARY =======================
\section{Summary}
AI engineer who builds AI-native developer tooling: coding agents, MCP integrations, and LLM workflows, owned from architecture to production. As an early-stage hire I shaped the AI/ML stack and shipped agents on the Anthropic API, including an MCP server exposing 180+ tools, with the evals and benchmarks to prove they work. I make pragmatic build-versus-integrate calls in ambiguous areas, prototype fast, iterate on feedback, and care about solving real user problems rather than building interesting technology for its own sake. Strong Python, some Kotlin/JVM. MSc Data Science (King's College London); 67K+ HuggingFace downloads across open models I have fine-tuned.

%======================= COMPETENCIES =======================
\section{Core Strengths}
\begin{tight}
  \item[\textcolor{accent}{\small$\bullet$}] \textbf{AI Developer Tooling:} coding agents \& assistants, AI-native workflows, MCP integrations, tool design, context management
  \item[\textcolor{accent}{\small$\bullet$}] \textbf{Agentic Systems \& LLM Pipelines:} multi-step agents (ReAct), prompt engineering, model selection, RAG, inference optimisation
  \item[\textcolor{accent}{\small$\bullet$}] \textbf{LLM Evaluation \& Benchmarking:} automated evals, LLM-as-judge, polyglot benchmarks, QoS metrics, observability
  \item[\textcolor{accent}{\small$\bullet$}] \textbf{Engineering \& Direction:} Python (TDD), Rust, some Kotlin/JVM; architecture decisions, influence without formal authority, product sense
\end{tight}

%======================= EXPERIENCE =======================
\section{Experience}
\entry{Built AI (FinTech \& AI Company)}{April 2022 -- Present}{AI/ML Engineer (Early Stage Hire)}{London, UK}
\begin{tight}
  \item Shaped and owned the AI/ML stack as an early-stage hire, from architecture and pragmatic build-versus-integrate decisions through to production for 5+ enterprise clients.
  \item Built and shipped agents and AI tooling on the Anthropic API, including an MCP server exposing 180+ tools through the Model Context Protocol for Claude Desktop and Glean-compatible clients.
  \item Designed LLM pipelines and context management (RAG, grounding, tool use) and evaluated them with automated evals, benchmarks, and LLM-as-judge scoring.
  \item Established evaluation and observability practices: Quality of Service metrics, OpenTelemetry to Arize Phoenix, and incident-triggering monitoring.
  \item Prototyped fast in ambiguous areas and iterated on feedback, optimising inference for latency and cost (local embeddings, p95 around 175ms) and improving a modelling engine's performance by 85\%.
\end{tight}

\entry{King's College London}{March 2021 -- January 2022}{Research Associate}{London, UK}
\begin{tight}
  \item Developed a deep learning model using Graph Convolution Networks and GANs for drug discovery, achieving a 63\% improvement in the novelty score of generated structures, over a preprocessing pipeline built for a million-row clinical dataset.
\end{tight}

%======================= PROJECTS =======================
\section{Open-Source \& Projects}
\begin{tight}
  \item \textbf{Spec-to-Code Agent:} autonomous ReAct-based coding agent (LangGraph + Anthropic API) that generates installable, tested packages from plain-language specs, benchmarked at a 9/9 pass rate across a polyglot suite (Python, Go, Rust, Java, JS). \href{https://github.com/abhijeetscode/Spec-to-Code-Agent}{View on GitHub}
  \item \textbf{Agentic Enterprise Assistant:} agent that answers questions over grounded data with tool use and context management, RBAC at an MCP server (Keycloak JWT), a grounding-validation gate, and auditable traces; OpenTelemetry to Arize Phoenix. Passed 13/13 evaluation cases. \href{https://github.com/abhijeetscode/Agentic-Enterprise-Assistant}{View on GitHub}
  \item \textbf{Transformer from Scratch in Rust:} GPT-style decoder-only transformer built from first principles (multi-head causal self-attention, pre-norm residual blocks, full training loop with SafeTensors checkpointing) using the Candle framework with Metal GPU acceleration. \href{https://github.com/abhijeetscode/transformer-rust}{View on GitHub}
\end{tight}

%======================= EDUCATION =======================
\section{Education}
\entry{King's College London}{}{MSc in Data Science}{London, UK}
\vspace{-3pt}
\entry{C-DAC, India}{}{Post Graduate Diploma in Advanced Computing (PG-DAC)}{Pune, India}
\vspace{-3pt}
\entry{University of Pune}{}{B.E. in Computer Science}{Pune, India}

\end{document}
